\section{Introduction}
\label{sec:introduction}

Large language models (LLMs) have become a general-purpose computing substrate for conversational assistants, code generation, retrieval-augmented generation, document and scientific-literature analysis, automated scientific discovery, and increasingly autonomous multi-agent workflows~\cite{ouyang2022instructgpt,chen2021codex,lewis2020rag,skarlinski2024paperqa,lu2024aiscientist,wu2024autogen}. These applications transform pretrained models into continuously accessed online services rather than one-time offline computations. A production service must accept requests with heterogeneous prompt and response lengths, maintain interactive latency under changing concurrency, and support models whose parameters and runtime states span multiple accelerators. As models acquire longer contexts and applications compose multiple model calls, inference efficiency has therefore become as important as model capability.

To support this demand, systems research has optimized LLM inference along several complementary dimensions. Scheduling and continuous batching improve accelerator utilization and balance throughput against latency~\cite{yu2022orca,agrawal2024sarathi,sun2024llumnix,duan2024muxserve,hong2025sola}. KV-cache research reduces fragmentation, reuses shared prefixes, separates cache storage from computation, or selectively retains important states to control the memory cost of long-context generation~\cite{kwon2023vllm,zheng2024sglang,qin2024mooncake,zhang2023h2o,xiao2024streamingllm,li2024snapkv,chen2024arkvale,lee2024infinigen,liu2024cachegen,lin2025qserve}. Multi-GPU parallelism partitions model tensors, routed experts, or long sequences to overcome the compute and memory limits of a single device~\cite{shoeybi2019megatron,rajbhandari2022deepspeedmoe,liu2024ringattention,wu2024loongserve,su2025seesaw}. Disaggregated architectures further separate inference phases or operators with different bottlenecks so that they can be provisioned independently~\cite{zhong2024distserve,zhu2025megascale,patel2024splitwise,strati2024dejavu,lin2024infinitellm}. Together, these advances improve throughput, latency, memory capacity, and scalability, making increasingly large and complex inference services practical. Their performance benefits, however, do not directly establish their energy efficiency.

The rapid expansion and continuous operation of LLM services create a growing energy challenge. Online inference repeatedly streams large weight tensors, reads and writes expanding KV caches, exchanges intermediate state across GPUs, and keeps provisioned accelerators powered while traffic fluctuates. Small inefficiencies at the request level accumulate across millions of invocations and long service lifetimes. The International Energy Agency estimates that data centers consumed about 415~TWh in 2024, approximately 1.5\% of global electricity use, and projects consumption to reach roughly 945~TWh by 2030, with AI as the most important driver of this growth~\cite{iea2025energyai}. Energy is thus becoming a first-order constraint on operating cost, power provisioning, and sustainable growth. Importantly, lowering instantaneous power does not guarantee lower energy: a slower operating point may extend request execution, and an architecture that improves utilization may introduce communication, synchronization, or idle-device overhead~\cite{zhang2026distributedenergy,ifath2026characterizing}.

Hardware diversity further complicates this trade-off. GPUs differ substantially in arithmetic throughput, memory capacity, memory bandwidth, interconnect support, and power-management interfaces. A compute-rich GPU can accelerate prefill or FFN computation but remain inefficient for memory-bound decoding if its capacity or bandwidth is insufficient; conversely, a memory-rich datacenter GPU can support large models and KV caches even when its peak arithmetic throughput is lower. Datacenter GPUs generally emphasize ECC memory, sustained operation, manageability, isolation, and scale-up/scale-out communication, whereas consumer GPUs prioritize high local throughput and cost but usually provide less memory and fewer datacenter interconnect and management features~\cite{nvidiaA800datasheet,nvidiaL40Sspec,nvidiaRTXBlackwell,nvidiaL20spec}. Frequency controls also differ across products rather than following a universal class rule. On our platforms, the A800 exposes practical runtime control of the SM/core clock while its memory clock remains fixed, whereas the RTX~5090 platform also permits memory-clock adjustment. These differences can change not only absolute performance and energy, but also which serving architecture and role placement are preferable.

Existing energy research can be grouped into three broad directions. Measurement studies quantify request-, token-, workflow-, or component-level energy and show that workload shape, batching, model structure, and hardware all matter~\cite{ifath2026characterizing,vellaisamy2026characterization,chung2025mlenergy}. Runtime-management systems adapt resource allocation, model parallelism, batching, power caps, or GPU frequency under service-level objectives~\cite{stojkovic2025dynamollm,kakolyris2025throttll,hankendi2026pals,jiang2025thunderserve,mei2025helix}. Phase- and operator-aware systems exploit the asymmetric behavior of prefill, decode, attention, and FFN computation through finer-grained placement and frequency control~\cite{liu2025greenllm,basit2026biscale,hu2026energy}. These studies establish important mechanisms and observations, but optimization systems generally target a particular architecture or policy, while measurement studies do not systematically compare alternative disaggregation boundaries. We still lack a controlled, architecture-level account of how different disaggregated designs behave across deployment plans, hardware platforms, communication environments, and runtime configurations. Without such an account, it remains unclear when specialization yields genuine energy savings and when it merely shifts cost among devices or converts computation into communication and idle time.

We address this gap through a progressive measurement study of non-disaggregated (Native), prefill--decode-disaggregated (PD), and attention--FFN-disaggregated (AF) serving in a common runtime. We begin with a broad characterization across dense and MoE models and representative input/output regimes to establish the performance--energy landscape and identify an architecture-sensitive MoE setting. We then examine how communication environments change the relative behavior of the architectures, followed by a study of data, tensor, expert, and hybrid parallel deployments. Finally, we move from software configuration to hardware characteristics, comparing GPUs with different compute and memory capabilities and exploring heterogeneous assignments in which specialized roles use different devices. This progression moves from application-visible variation to communication, deployment, and hardware causes, while matched traces, SLO-aware filtering, repeated open-loop runs, and deployment-level energy accounting keep comparisons controlled.

Our study makes the following contributions:
\begin{itemize}
  \item \textbf{An architecture-centered energy characterization.} We formulate and implement a common measurement boundary for Native, PD, and AF, enabling direct comparison of complete deployments rather than isolated kernels or worker pools. This exposes the fundamental trade-off between resource specialization and the transfer, synchronization, and device-residency costs introduced by disaggregation.
  \item \textbf{Evidence that architecture choice is workload dependent.} Completed repeated RDMA measurements exhibit a crossover: PD consumes less GPU energy for short QA and decode-heavy chatbot traffic, whereas AF consumes less for prefill-heavy RAG, long-output summarization, and long-context traffic; the two are nearly tied on the balanced workload. Thus, neither form of disaggregation is uniformly preferable, and input/output structure must be part of an energy-aware architecture decision.
  \item \textbf{A cross-layer explanation of disaggregated-serving energy.} By progressively controlling network conditions and parallel placement, our study separates useful phase/operator specialization from state transfer, collective communication, pipeline bubbles, and idle-power overhead. This analysis connects request-level energy outcomes to their underlying mechanisms rather than reporting aggregate joules alone.
  \item \textbf{A hardware-aware path toward heterogeneous serving.} We relate architecture behavior to GPU compute throughput, memory capacity, bandwidth, interconnect, and controllable frequency domains. The resulting analysis provides an empirical basis for assigning different GPU classes to specialized roles and for future architecture- and hardware-aware controllers.
\end{itemize}
% Before final submission, augment the last two contribution bullets with the
% quantitative RQ2--RQ5 findings once those experiment groups are complete.

The remainder of this paper is organized as follows. Section~2 introduces LLM inference, parallel deployment, disaggregated serving, energy management, and GPU classes. Section~3 reviews related work and positions our measurement scope. Section~4 describes the hardware, models, workloads, architectures, and measurement methodology. Section~5 presents the experimental questions and results. Section~6 discusses their implications, limitations, and future directions, and Section~7 concludes.
